\documentclass[a4paper,12pt,twoside]{memoir}

% Castellano
\usepackage[spanish,es-tabla]{babel}
\selectlanguage{spanish}
\usepackage[utf8]{inputenc}
\usepackage[T1]{fontenc}
\usepackage{lmodern} % scalable font
\usepackage{microtype}
\usepackage{placeins}
\usepackage{float}

\RequirePackage{booktabs}
\RequirePackage[table]{xcolor}
\RequirePackage{xtab}
\RequirePackage{multirow}

% Links
\PassOptionsToPackage{hyphens}{url}\usepackage[colorlinks]{hyperref}
\hypersetup{
	allcolors = {red}
}

% Ecuaciones
\usepackage{amsmath}

% Rutas de fichero / paquete
\newcommand{\ruta}[1]{{\sffamily #1}}

% Párrafos
\nonzeroparskip

% Huérfanas y viudas
\widowpenalty100000
\clubpenalty100000

% Evitar solapes en el header
\nouppercaseheads

% Imagenes
\usepackage{graphicx}
\newcommand{\imagen}[3]{
	\begin{figure}[!ht]
		\centering
		\includegraphics[width=#3\textwidth]{#1}
		\caption{#2}\label{fig:#1}
	\end{figure}
	\FloatBarrier
}

\newcommand{\imagenflotante}[2]{
	\begin{figure}
		\centering
		\includegraphics[width=0.9\textwidth]{#1}
		\caption{#2}\label{fig:#1}
	\end{figure}
}



% El comando \figura nos permite insertar figuras comodamente, y utilizando
% siempre el mismo formato. Los parametros son:
% 1 -> Porcentaje del ancho de página que ocupará la figura (de 0 a 1)
% 2 --> Fichero de la imagen
% 3 --> Texto a pie de imagen
% 4 --> Etiqueta (label) para referencias
% 5 --> Opciones que queramos pasarle al \includegraphics
% 6 --> Opciones de posicionamiento a pasarle a \begin{figure}
\newcommand{\figuraConPosicion}[6]{%
  \setlength{\anchoFloat}{#1\textwidth}%
  \addtolength{\anchoFloat}{-4\fboxsep}%
  \setlength{\anchoFigura}{\anchoFloat}%
  \begin{figure}[#6]
    \begin{center}%
      \Ovalbox{%
        \begin{minipage}{\anchoFloat}%
          \begin{center}%
            \includegraphics[width=\anchoFigura,#5]{#2}%
            \caption{#3}%
            \label{#4}%
          \end{center}%
        \end{minipage}
      }%
    \end{center}%
  \end{figure}%
}

%
% Comando para incluir imágenes en formato apaisado (sin marco).
\newcommand{\figuraApaisadaSinMarco}[5]{%
  \begin{figure}%
    \begin{center}%
    \includegraphics[angle=90,height=#1\textheight,#5]{#2}%
    \caption{#3}%
    \label{#4}%
    \end{center}%
  \end{figure}%
}
% Para las tablas
\newcommand{\otoprule}{\midrule [\heavyrulewidth]}
%
% Nuevo comando para tablas pequeñas (menos de una página).
\newcommand{\tablaSmall}[5]{%
 \begin{table}[!ht]
  \begin{center}
   \rowcolors {2}{gray!35}{}
   \begin{tabular}{|m{4cm} | m{8cm}|}
    \toprule
    #4
    \otoprule
    #5
    \bottomrule
   \end{tabular}
   \caption{#1}
   \label{tabla:#3}
  \end{center}
 \end{table}
}

%
%Para el float H de tablaSmallSinColores
\usepackage{float}

%
% Nuevo comando para tablas pequeñas (menos de una página).
\newcommand{\tablaSmallSinColores}[5]{%
 \begin{table}[H]
  \begin{center}
   \begin{tabular}{#2}
    \toprule
    #4
    \otoprule
    #5
    \bottomrule
   \end{tabular}
   \caption{#1}
   \label{tabla:#3}
  \end{center}
 \end{table}
}

\newcommand{\tablaApaisadaSmall}[5]{%
\begin{landscape}
  \begin{table}
   \begin{center}
    \rowcolors {2}{gray!35}{}
    \begin{tabular}{#2}
     \toprule
     #4
     \otoprule
     #5
     \bottomrule
    \end{tabular}
    \caption{#1}
    \label{tabla:#3}
   \end{center}
  \end{table}
\end{landscape}
}

%
% Nuevo comando para tablas grandes con cabecera y filas alternas coloreadas en gris.
\newcommand{\tabla}[6]{%
  \begin{center}[H]
    \tablefirsthead{
      \toprule
      #5
      \otoprule
    }
    \tablehead{
      \multicolumn{#3}{l}{\small\sl continúa desde la página anterior}\\
      \toprule
      #5
      \otoprule
    }
    \tabletail{
      \hline
      \multicolumn{#3}{r}{\small\sl continúa en la página siguiente}\\
    }
    \tablelasttail{
      \hline
    }
    \bottomcaption{#1}
    \rowcolors {2}{gray!35}{}
    \begin{xtabular}{#2}
      #6
      \bottomrule
    \end{xtabular}
    \label{tabla:#4}
  \end{center}
}

%
% Nuevo comando para tablas grandes con cabecera.
\newcommand{\tablaSinColores}[6]{%
  \begin{center}
    \tablefirsthead{
      \toprule
      #5
      \otoprule
    }
    \tablehead{
      \multicolumn{#3}{l}{\small\sl continúa desde la página anterior}\\
      \toprule
      #5
      \otoprule
    }
    \tabletail{
      \hline
      \multicolumn{#3}{r}{\small\sl continúa en la página siguiente}\\
    }
    \tablelasttail{
      \hline
    }
    \bottomcaption{#1}
    \begin{xtabular}{#2}
      #6
      \bottomrule
    \end{xtabular}
    \label{tabla:#4}
  \end{center}
}

%
% Nuevo comando para tablas grandes sin cabecera.
\newcommand{\tablaSinCabecera}[5]{%
  \begin{center}
    \tablefirsthead{
      \toprule
    }
    \tablehead{
      \multicolumn{#3}{l}{\small\sl continúa desde la página anterior}\\
      \hline
    }
    \tabletail{
      \hline
      \multicolumn{#3}{r}{\small\sl continúa en la página siguiente}\\
    }
    \tablelasttail{
      \hline
    }
    \bottomcaption{#1}
  \begin{xtabular}{#2}
    #5
   \bottomrule
  \end{xtabular}
  \label{tabla:#4}
  \end{center}
}



\definecolor{cgoLight}{HTML}{EEEEEE}
\definecolor{cgoExtralight}{HTML}{FFFFFF}

%
% Nuevo comando para tablas grandes sin cabecera.
\newcommand{\tablaSinCabeceraConBandas}[5]{%
  \begin{center}[H]
    \tablefirsthead{
      \toprule
    }
    \tablehead{
      \multicolumn{#3}{l}{\small\sl continúa desde la página anterior}\\
      \hline
    }
    \tabletail{
      \hline
      \multicolumn{#3}{r}{\small\sl continúa en la página siguiente}\\
    }
    \tablelasttail{
      \hline
    }
    \bottomcaption{#1}
    \rowcolors[]{1}{cgoExtralight}{cgoLight}

  \begin{xtabular}{#2}
    #5
   \bottomrule
  \end{xtabular}
  \label{tabla:#4}
  \end{center}
}




\graphicspath{ {./img/} }

% Capítulos
\chapterstyle{bianchi}
\newcommand{\capitulo}[2]{
	\setcounter{chapter}{#1}
	\setcounter{section}{0}
	\setcounter{figure}{0}
	\setcounter{table}{0}
	\chapter*{#2}
	\addcontentsline{toc}{chapter}{#2}
	\markboth{#2}{#2}
}

% Apéndices
\renewcommand{\appendixname}{Apéndice}
\renewcommand*\cftappendixname{\appendixname}

\newcommand{\apendice}[1]{
	%\renewcommand{\thechapter}{A}
	\chapter{#1}
}

\renewcommand*\cftappendixname{\appendixname\ }

% Formato de portada
\makeatletter
\usepackage{xcolor}
\newcommand{\tutor}[1]{\def\@tutor{#1}}
\newcommand{\course}[1]{\def\@course{#1}}
\definecolor{cpardoBox}{HTML}{E6E6FF}
\def\maketitle{
  \null
  \thispagestyle{empty}
  % Cabecera ----------------
\noindent\includegraphics[width=\textwidth]{cabecera}\vspace{1cm}%
  \vfill
  % Título proyecto y escudo informática ----------------
  \colorbox{cpardoBox}{%
    \begin{minipage}{.8\textwidth}
      \vspace{.5cm}\Large
      \begin{center}
      \textbf{TFG del Grado en Ingeniería Informática}\vspace{.6cm}\\
      \textbf{\LARGE\@title{}}
      \end{center}
      \vspace{.2cm}
    \end{minipage}

  }%
  \hfill\begin{minipage}{.20\textwidth}
    \includegraphics[width=\textwidth]{escudoInfor}
  \end{minipage}
  \vfill
  % Datos de alumno, curso y tutores ------------------
  \begin{center}%
  {%
    \noindent\LARGE
    Presentado por \@author{}\\ 
    en Universidad de Burgos --- \@date{}\\
    Tutores: \@tutor{}\\
  }%
  \end{center}%
  \null
  \cleardoublepage
  }
\makeatother

\newcommand{\nombre}{Diego Garrido Calvo} %%% cambio de comando

% Datos de portada
\title{Control de un molino eólico utilizando aprendizaje por refuerzo profundo \\Documentación Técnica}
\author{\nombre}
\tutor{Daniel Sarabia Ortiz y Roberto Carlos Casado Vara}
\date{\today}

\begin{document}

\maketitle



\cleardoublepage



%%%%%%%%%%%%%%%%%%%%%%%%%%%%%%%%%%%%%%%%%%%%%%%%%%%%%%%%%%%%%%%%%%%%%%%%%%%%%%%%%%%%%%%%



\frontmatter


\clearpage

% Indices
\tableofcontents

\clearpage

\listoffigures

\clearpage

\listoftables

\clearpage

\mainmatter

\appendix

\apendice{Plan de Proyecto Software}

\section{Introducción}

En este apéndice, se van a explicar y comentar todos los aspectos relacionados con la planificación del proyecto desde el punto de vista del desarrollo software.

La fase de planificación se divide en:
\begin{itemize}
    \item Planificación temporal.
    \item Estudio de viabilidad.
\end{itemize}

La primera parte está destinada, mayormente, a la distribución del proyecto en cuanto al tiempo, la división en tareas del mismo y la organización tanto del alumno como de los tutores para su elaboración.
La segunda parte, se centra en el estudio de la rentabilidad del proyecto. Esta, se divide en dos puntos importantes:
\begin{itemize}
    \item Viabilidad económica: en esta sección se hace un baremo entre los costes y los beneficios del proyecto.
    \item Viabilidad legal: está mayormente orientada a temas de licencias y protección de datos.
\end{itemize}


\section{Planificación temporal}

En cuanto a la organización del proyecto, al comienzo del mismo se planteó utilizar la metodología \textit{Scrum} para organizar las tareas y la regularidad de las reuniones. 
Se ha procurado seguir las pautas establecidas sobre la organización de los \textit{sprints} en todo momento, aunque también es cierto que, debido a los problemas que han ido surgiendo, ha sido imposible seguir esto a rajatabla.

\begin{itemize}
    \item Estrategia de desarrollo incremental a través de \textit{sprints}.
    \item Sprints cada 1-2 semanas según la etapa del proyecto.
    \item En cada \textit{sprints}, se realizaba una entrega de una parte del proyecto funcional.
    \item Se han realizado reuniones en cada uno de los \textit{sprints}.
    \item En cada uno de los \textit{sprints} se han generado una serie de tareas a realizar.
\end{itemize}

La estimación se realizó mediante una serie de \textit{story points}\cite{storyPoints}, cuya traducción temporal aparece en la tabla \ref{tabla:storypoint}.

\tablaSmall{Story Point Cheat Sheet}{l c}{storypoint}
{\multicolumn{1}{l}{Story points} & Estimación temporal\\}{ 
1 & 30 minutos\\
3 & 2 horas\\
5 & 3 horas\\
8 & 5 horas\\
13 & 10 horas\\
20 & 20 horas\\
40 & 32 horas\\
} 

Los \textit{sprints} realizados han sido los siguientes:

\subsection{Sprint 0 (15/02/2022 - 28/02/2022)}
\begin{itemize}
    \item Elección de la temática del proyecto entre las ofertadas por la UBU.
    \item Primera reunión con los tutores, para comentar las bases generales del TFG, plazos de entrega, división del proyecto en \textit{sprints}.
    \item Firma, por parte de los tutores, del documento que acreditaba que el alumno había sido asignado a este proyecto.
\end{itemize}

\subsection{Sprint 1 (01/03/2022 - 15/03/2022)}
\begin{itemize}
    \item Investigación exhaustiva sobre el tema del proyecto.
    \item Lectura de libros y artículos recomendados por los tutores para reforzar conocimientos.
    \item Pequeño seminario en el cuál se aclararon dudas y se explicaron de forma más concreta ciertos puntos más difíciles de comprender.
    \item Creación del repositorio de GitHub.
    \item Creación del Sharepoint.
    \item Descarga e instalación de las herramientas y el entorno de desarrollo.
    \item Establecimiento de los objetivos y metas del proyecto.
\end{itemize}

\imagen{img_story_points/StoryPoint2.png}{Gráfica Burndown Sprint 1.}{0.8}

\subsection{Sprint 2 (15/03/2022 - 31/03/2022)}
\begin{itemize}
    \item instalación de paquetes y librerías necesarias y comienzo de la parte de programación del proyecto.
    \item Investigación sobre el uso del editor de textos \LaTeX y primer contacto con la herramienta.
    \item Lectura e investigación sobre distintos Trabajos de Fin de Grado disponibles.
    \item Comienzo de la memoria con los Conceptos teóricos que se aprendieron hasta el momento y las Técnicas y herramientas que se iban usando (Estos dos apartados de la memoria se han ido completando a lo largo del proyecto según se aprendían conceptos y se usaban nuevas herramientas).
\end{itemize}

\imagen{img_story_points/StoryPoint3.png}{Gráfica Burndown Sprint 2.}{0.8}

\subsection{Sprint 3 (01/04/2022 - 15/04/2022)}
\begin{itemize}
    \item Finalización del desarrollo del entorno de la parte de aprendizaje por refuerzo.
    \item Investigación sobre cómo construir la red neuronal y primeros pasos del desarrollo.
\end{itemize}

\imagen{img_story_points/StoryPoint4.png}{Gráfica Burndown Sprint 3.}{0.8}

\subsection{Sprint 4 (16/04/2022 - 30/04/2022)}
\begin{itemize}
    \item Creación tanto de la red neuronal como del agente a entrenar.
    \item Primeras pruebas de la calidad del entrenamiento sobre el modelo.
    \item Contrastación de las pruebas con los tutores, para recibir soluciones de errores y sugerencias de mejoras.
    \item Continuación de la memoria del proyecto.
\end{itemize}

\imagen{img_story_points/StoryPoint5.png}{Gráfica Burndown Sprint 4.}{0.8}

\subsection{Sprint 5 (01/05/2022 - 15/05/2022)}
\begin{itemize}
    \item Finalización de la funcionalidad de la parte de aprendizaje por refuerzo, a falta de alguna mejora.
    \item Investigación sobre el funcionamiento de una turbina eólica tanto teórica, como matemáticamente.
    \item Contrastación de lo desarrollado hasta el momento con otros ejemplos de entrenamiento por refuerzo disponibles en internet.
    \item Investigación sobre el método de Euler.
\end{itemize}

\imagen{img_story_points/StoryPoint6.png}{Gráfica Burndown Sprint 5.}{0.8}

\subsection{Sprint 6 (16/05/2022 - 31/05/2022)}
\begin{itemize}
    \item Implementación de la matemática de la turbina eólica en el código.
    \item Realización de varias pruebas para comprobar la funcionalidad del algoritmo.
    \item Implementación de algunas mejoras sugeridas en la reunión Scrum.
    \item Comienzo de desarrollo de los anexos.
    \item Desarrollo de la interfaz web.
    \item Creación de lista con detalles y mejoras a falta de implementar en el proyecto.
\end{itemize}

\imagen{img_story_points/StoryPoint7.png}{Gráfica Burndown Sprint 6.}{0.8}

\subsection{Sprint 7 (01/06/2022 - 15/06/2022)}
\begin{itemize}
    \item Implementación de mejoras restantes.
    \item Resolución de ciertas dudas sobre la memoria y algunos puntos de las bases de la realización del TFG pendientes de aclarar.
    \item Continuación tanto de la memoria, como de los anexos.
\end{itemize}

\imagen{img_story_points/StoryPoint8.png}{Gráfica Burndown Sprint 7.}{0.8}

\subsection{Sprint 8 (15/06/2022 - 31/06/2022)}
\begin{itemize}
    \item Finalización y corrección de la memoria.
    \item Creación del documento para la presentación ante el tribunal.
\end{itemize}

\imagen{img_story_points/StoryPoint9.png}{Gráfica Burndown Sprint 8.}{0.8}

\section{Estudio de viabilidad}

A continuación, se realiza un análisis económico en el hipotético caso de que el proyecto se llegase a utilizar en el mercado laboral. 

Aún así, dejar claro que para ello sería necesario de 1 a 2 años más de desarrollo del proyecto para que fuese viable.

\subsection{Viabilidad económica}

En este apartado de la viabilidad, se analizan tanto los costes como los beneficios que podría aportar el proyecto en el caso de ser usado en un entorno empresarial real.

\subsubsection{Coste de software}

El coste de software se refiere al precio de las licencias de los programas, aplicaciones y suscripciones web que se han utilizado durante el desarrollo del proyecto. Los costes de software de nuestro proyecto son los que aparecen en la tabla \ref{tabla:costesoftware}

\tablaSmall{Coste de software}{l c}{costesoftware}
{ \multicolumn{1}{l}{Concepto} & Coste\\}{ 
Microsoft Windows 10 Pro & 11.50 €\\
Codacy & 14.40 €\\
\textbf{Total} & \textbf{25.90 €}\\
} 

\subsubsection{Coste de personal}

El coste de personal se refiere al salario que habría que pagarle al desarrollador junior por las horas imputadas. Para proyecto han sido necesarias 30 horas semanales durante 4 meses. En la tabla \ref{tabla:costepersonal} podemos ver una estimación de dichos costes.

\tablaSmall{Coste de personal}{l c}{costepersonal}
{ \multicolumn{1}{l}{Concepto} & Pagos\\}{ 
Salario neto del empleado & 1000.00 €\\
IRPF (19\%) & 360.00 €\\
Seguridad social (28,5\%) & 537.00 €\\
Salario bruto del trabajador & 1897.00 €\\
\textbf{Total} & \textbf{7588.00 €}\\
}

\subsubsection{Coste de material}

El coste de material son los productos de hardware que han sido necesarios para la realización del proyecto. En este caso, como se puede ver en la tabla \ref{tabla:costematerial}, el único gasto en material sería el precio del ordenador utilizado para el desarrollo.

\tablaSmall{Coste de material}{l c}{costematerial}
{ \multicolumn{1}{l}{Concepto} & Coste\\}{ 
Ordenador & 700.00 €\\
\textbf{Total} & \textbf{700.00 €}\\
} 

\subsubsection{Beneficios}

Debido a que este es un simple proyecto de fin de grado destinado al aprendizaje y a poner a prueba al alumno, no se busca obtener ningún beneficio de ello.

\subsection{Viabilidad legal}

A continuación, se pasa a analizar las leyes que podrían afectar al proyecto desarrollado, en caso de utilizarse a nivel laboral.

\subsubsection{Licencias}

Uno de los mayores problemas que suelen surgir con este tipo de proyectos, es el tema de las licencias de software.
En este caso, casi todo el software utilizado es de uso libre y, en el caso de haberse utilizado algún software de pago, todas las licencias están en orden, por lo que no supondría un problema legal en este caso.

\subsubsection{Protección de datos}

Otro tema que también es bastante relevante en cuanto a este apartado, es el tema de la protección de datos personales.
En este proyecto, no se solicita al usuario de ningún dato personal, por lo que tampoco habría que preocuparse en ese aspecto.

\apendice{Especificación de Requisitos}

\section{Introducción}

En esta parte de los anexos se detalla toda la serie de requisitos que definen el proyecto. Sirven para contrastar cierta documentación entre el cliente y el equipo de desarrollo y para analizar el comportamiento de la aplicación.

Para realizar una buena especificación de requisitos, esta debe ser:
\begin{itemize}
    \item \textbf{Completa}: deben aparecer todos los requisitos, tanto funcionales como no funcionales.
    \item \textbf{Correcta}: el software del proyecto debe cumplir los requisitos que aquí se exponen.
    \item \textbf{Consistente}: debe seguir una línea de coherencia y cohesión con los propios requerimientos.
    \item \textbf{Modificable}: facilidad de modificación de los requerimientos.
    \item \textbf{Verificable}: debe ser posible probarlo sin costo alguno.
    \item \textbf{Trazable}: posibilidad de conocer la ubicación de un ítem a través de su identificador.
    \item \textbf{Priorizable}: los requisitos se ordenan según su importancia para dar preferencia a los más relevantes.
\end{itemize}


\section{Objetivos generales}

Los objetivos del proyecto son los siguientes:
\begin{itemize}
    \item Modelar el comportamiento de una turbina eólica a partir de un algoritmo de aprendizaje por refuerzo, el cuál sea capaz de ajustarse automáticamente a una demanda de energía dada.
    \item Calcular y aplicar las ecuaciones que intervienen en el funcionamiento de una turbina eólica.
    \item Diseñar una interfaz-web sencilla para interactuar con el software del modelo.
    \item Facilitar el entendimiento de la respuesta del molino mediante gráficas.
\end{itemize}


\section{Catálogo de requisitos}

Se proceden a describir los requisitos del proyecto, los cuales se dividen en dos tipos:

\subsection{Requisitos funcionales}

\begin{itemize}
    \item \textbf{RF 1 - Entrenamiento del agente}. El agente debe ser capaz de entrenarse sin problemas, devolviendo un buen resultado de entrenamiento.
    \begin{itemize}
        \item \textbf{RF 1.1 - Creación del entorno}. El entorno debe crearse correctamente, instanciando la clase en una variable.
        \item \textbf{RF 1.2 - Creación del agente}. El agente debe comunicarse correctamente con la red neuronal con la que va a trabajar.
        \item \textbf{RF 1.3 - Entrenamiento}. El agente debe entrenarse, comunicándose correctamente con el entorno donde va a realizar las pruebas del entrenamiento. También es importante el tiempo que va a entrenar, ya que demasiado tiempo de entrenamiento puede llevar a lo que conocemos como \textit{overtraining}\cite{overtraining}.
    \end{itemize}
    \item \textbf{RF 2 - Pruebas del entrenamiento}. Deben realizarse correctamente una o varias pruebas que muestren el correcto entrenamiento del agente.
    \begin{itemize}
        \item \textbf{RF 2.1 - Pruebas pre-entrenamiento}. Se realizan unas pruebas tomando acciones aleatorias para mostrar el funcionamiento del modelo sin un agente entrenado.
        \item \textbf{RF 2.2 - Pruebas post-entrenamiento}. El agente entrenado debe realizar una serie de pruebas sobre el modelo para mostrar los resultados tras el entrenamiento.
    \end{itemize}
    \item \textbf{RF 3 - Correcto funcionamiento de la interfaz web}. El usuario debe ser capaz de introducir una potencia de referencia y observar en una gráfica la respuesta del modelo ante esa acción.
    \begin{itemize}
        \item \textbf{RF 3.1 - Comprobación de seguridad} Se comprueba que el usuario pueda introducir únicamente una variable numérica.
        \item \textbf{RF 3.2 - Impresión por pantalla de gráfica} Se procede a mostrar correctamente una o varias gráficas en la propia web con resultados sobre la ejecución del modelo.
    \end{itemize}
\end{itemize}

\subsection{Requisitos no funcionales}

\begin{itemize}
    \item  \textbf{RNF 1 - Usabilidad}. La interfaz de la aplicación debe ser sencilla e intuitiva.
    \item  \textbf{RNF 2 - Rendimiento}. Los tiempos de carga deben ser aceptables y no tediosos.
    \item  \textbf{RNF 3 - Disponibilidad}. La aplicación debe poder usarse en cualquier lugar, en cualquier momento.
    \item  \textbf{RNF 4 - Accesibilidad}. La aplicación debe poder usarse por cualquier persona independientemente de sus capacidades técnicas/cognitivas. 
\end{itemize}


\section{Especificación de requisitos}

A continuación, se procede a detallar los distintos casos de uso que cubren los requisitos funcionales que acabamos de indicar.

\tablaSmall{CU 01 - Creación del entorno}{l c}{creacionentorno}
{ \multicolumn{1}{l}{CU 01} & Creación del entorno\\}{ 
Autor & Diego Garrido Calvo\\
Requisitos asociados & RF 1.1\\
Descripción & El entorno debe crearse correctamente, instanciando la clase en una variable.\\
Pre-condición & El usuario debe haber ejecutado el modelo.\\
Acción & El usuario ejecuta el código de Jupyter Notebook.\\
Post-condición & El entorno queda instanciado para su posterior uso.\\
}

\tablaSmall{CU 02 - Creación del agente}{l c}{creaciónagente}
{ \multicolumn{1}{l}{CU 02} & Creación del agente\\}{ 
Autor & Diego Garrido Calvo\\
Requisitos asociados & RF 1.2\\
Descripción & El agente debe comunicarse correctamente con la red neuronal con la que va a trabajar.\\
Pre-condición & La red neuronal debe estar correctamente creada.\\
Acción & Se instancian tanto el agente como la red neuronal.\\
Post-condición & El agente queda listo para su entrenamiento.\\
} 

\tablaSmall{CU 03 - Entrenamiento}{l c}{entrenamiento}
{ \multicolumn{1}{l}{CU 03} & Entrenamiento\\}{ 
Autor & Diego Garrido Calvo\\
Requisitos asociados & RF 1.3\\
Descripción & El agente debe entrenarse, comunicándose correctamente con el entorno donde va a realizar las pruebas del entrenamiento. También es importante el tiempo que va a entrenar, ya que demasiado tiempo de entrenamiento puede llevar a lo que conocemos como \textit{overtraining}.\\
Pre-condición & Tanto el agente como el entorno deben estar creados.\\
Acción & El agente se entrena sobre el entorno instanciado.\\
Post-condición & El agente queda entrenado.\\
}

\tablaSmall{CU 04 - Pruebas pre-entrenamiento}{l c}{pruebaspre}
{ \multicolumn{1}{l}{CU 04} & Pruebas pre-entrenamiento\\}{ 
Autor & Diego Garrido Calvo\\
Requisitos asociados & RF 2.1\\
Descripción & Se realizan unas pruebas tomando acciones aleatorias para mostrar el funcionamiento del modelo sin un agente entrenado.\\
Pre-condición & Debe haberse creado el entorno.\\
Acción & El usuario elige la opción de ver el resultado antes del entrenamiento.\\
Post-condición & Resultados de la prueba.\\
}

\tablaSmall{CU 05 - Pruebas post-entrenamiento}{l c}{pruebaspost}
{ \multicolumn{1}{l}{CU 05} & Pruebas post-entrenamiento\\}{ 
Autor & Diego Garrido Calvo\\
Requisitos asociados & RF 2.2\\
Descripción & El agente entrenado debe realizar una serie de pruebas sobre el modelo para mostrar los resultados tras el entrenamiento.\\
Pre-condición & Deben haberse creado tanto el entorno como el agente, y el agente estar entrenado.\\
Acción & El usuario elige la opción de ver el resultado después del entrenamiento.\\
Post-condición & Resultados de la prueba.\\
}

\tablaSmall{CU 06 - Comprobación de seguridad}{l c}{comprobacionseguridad}
{ \multicolumn{1}{l}{CU 06} & Comprobación de seguridad\\}{ 
Autor & Diego Garrido Calvo\\
Requisitos asociados & RF 3.1\\
Descripción & Se comprueba que el usuario pueda introducir únicamente una variable numérica.\\
Pre-condición & El usuario debe haber ejecutado el modelo tras introducir una cantidad de potencia.\\
Acción & Se comprueba el tipo de la variable introducida.\\
Post-condición & En caso de que sea una variable numérica, el modelo se ejecuta; en caso contrario, aparece por pantalla un mensaje de error indicando al usuario lo que ha hecho mal.\\
}

\tablaSmall{CU 07 - Impresión por pantalla de gráfica}{l c}{impresionpantalla}
{ \multicolumn{1}{l}{CU 07} & Impresión por pantalla de gráfica\\}{ 
Autor & Diego Garrido Calvo\\
Requisitos asociados & RF 3.2\\
Descripción & Se procede a mostrar correctamente una o varias gráficas en la propia web con resultados sobre la ejecución del modelo.\\
Pre-condición & La comprobación de errores debe haber finalizado sin errores.\\
Acción & Se instancia un \textit{plot} con los datos obtenidos de la prueba.\\
Post-condición & Se muestra por pantalla la gráfica resultante.\\
}
\apendice{Especificación de diseño}

\section{Introducción}

Mediante la especificación de diseño, se aporta un mayor entendimiento del proyecto al equipo de desarrollo y se definen los datos que usa la aplicación, el diseño procedimental y la arquitectura.

\section{Diseño de datos}

En el caso de este proyecto, no es necesario el manejo de datos.

Como mucho podríamos hablar de la memoria secuencial que utiliza el agente. En esta se van almacenando los resultados de los entrenamientos pasados, hasta un límite que definamos nosotros.
Cuando el \textit{backlog} de los entrenamientos llega a este límite, la información del último entrenamiento sustituirá a la más antigua en memoria, ya que presupone que esta nueva información es más útil y precisa.

\section{Diseño procedimental}

En este apartado se procede a detallar la secuencia que sigue la ejecución del modelo. 

En el diagrama de secuencia\cite{diagramaSecuencia} de la figura \ref{fig:DiagramaSecuencia}, se aprecian las interacciones entre los objetos que conforman el modelo.

\imagen{DiagramaSecuenciaTFG.png}{Diagrama de secuencia del modelo}{1.0}
\label{fig:DiagramaSecuencia}

A continuación se pasa a explicar este diagrama con más detalle:

\begin{itemize}
    \item Lo primero que se hace es ejecutar la aplicación desde Jupyter Notebook. De esta forma se instancia el entorno, la red neuronal, el agente y se abre el servidor local donde se va a ejecutar la interfaz web.
    \item A continuación, el usuario abre la interfaz, introduce la potencia a generar y ejecuta uno de los dos botones disponibles. Independientemente del botón pulsado, se hace una comprobación de seguridad, para asegurar que no se haya introducido un valor no numérico o fuera del rango establecido, en caso de que el valor introducido no sea válido, aparecerá un mensaje de error \ref{fig:comprobación}.
    \item En caso de introducir un valor de potencia válido, dependerá del botón seleccionado el tipo de ejecución.
\end{itemize}

\imagen{img_capturas_interfaz/ErrorCheck.png}{Comprobación de seguridad}{0.7}
\label{fig:comprobación}

\subsection{Modelo desentrenado}

Si se elige ejecutar el modelo sin entrenar, lo primero que se hace es modificar el entorno para que la potencia de referencia sea la elegida por el usuario.
Tras esto se realiza una ejecución del modelo tomando acciones totalmente aleatorias sobre el ángulo del pitch y se muestra una gráfica con un resumen de la ejecución en 2000 pasos, o los pasos ejecutados en caso de haber alcanzado por casualidad la potencia de referencia antes de esos 2000.

\subsection{Modelo entrenado}

En el caso del modelo entrenado, también se modifica el entorno con la potencia de referencia que el usuario introduce. La diferencia es que en este caso, se entrena el agente creado anteriormente durante 20000 pasos (2 minutos aproximadamente).
Una vez entrenado el agente se vuelve a ejecutar el modelo, pero esta vez en lugar de tomar acciones aleatorias, es el agente el que elige que acción tomar en cada paso.

Por último, una vez terminada la ejecución del modelo, tanto para el modelo entrenado como para el modelo desentrenado, se muestra por pantalla una gráfica y una serie de datos con los resultados de la ejecución.

\section{Diseño de interfaces}

Para este proyecto, se pedía la creación de una interfaz sencilla que permitiese al usuario introducir una cantidad de potencia.

En la figura \ref{fig:interfaz}, se muestra otro diagrama del planteamiento de la interfaz.

\imagen{DiagramaInterfaz.png}{Diagrama de la interfaz de la aplicación}{1}
\label{fig:interfaz}

La interfaz se resume en un cuadro de texto donde se introducirá la potencia y dos botones. 
El botón \textbf{Sin entrenar} muestra una gráfica del funcionamiento del modelo tomando acciones aleatorias, mientras que el botón \textbf{Entrenado}, entrena el agente para esa potencia específica y muestra otra gráfica con el resultado tras el entrenamiento \ref{fig:ResultadoInterfaz}.

\imagen{img_capturas_interfaz/Resultado-interfaz.png}{Ejemplo de respuesta gráfica de la interfaz ante una ejecución sin entrenar.}{1}
\label{fig:ResultadoInterfaz}

En la gráfica de la izquierda esta representada la potencia calculada con una línea azul y la potencia de referencia con una línea roja. En la gráfica de la izquierda esta representado el valor del ángulo durante toda la ejecución, todo ello para un viento de 10 m/s.

\imagen{img_capturas_interfaz/entrenado-grafica.png}{Ejemplo de respuesta gráfica de la interfaz ante una ejecución con el modelo entrenado.}{1}
\label{fig:entrenadoGrafica}

Aquí \ref{fig:entrenadoGrafica} otro ejemplo de respuesta de la interfaz, pero esta vez con el modelo entrenado para una potencia de referencia de 2000 kW. Como se puede comprobar, esta vez la potencia generada alcanza la de referencia antes de los 2000 pasos.

Además de un resultado gráfico, también se aportan tanto la potencia inicial, como la potencia final y el ángulo del pitch óptimo de esa ejecución \ref{fig:datosInterfaz}.

\imagen{img_capturas_interfaz/DatosInterfaz.png}{Ejemplo de resultados para una ejecución con potencia de referencia 2000 kW.}{0.5}
\label{fig:datosInterfaz}

Todo ello con la intención de aportar al usuario datos más precisos, ya que a través de la gráfica solo podemos estimar el resultado final debido a que no se ve con claridad.


\apendice{Documentación técnica de programación}

\section{Introducción}

En este apéndice, se procede a detallar la documentación técnica de la programación del modelo.

Aquí se incluyen la instalación del entorno de desarrollo, así como la compilación y ejecución del proyecto, la estructura de directorios, el manual del programador y las pruebas realizadas.

\section{Estructura de directorios}

El repositorio cuenta con la siguiente distribución:
\begin{itemize}
    \item \textbf{/}: es el directorio raíz, donde están almacenados todos los archivos del proyecto junto al fichero \textit{README}.
    \item \textbf{/docs}: aquí podemos encontrar toda la documentación del proyecto (memoria y anexos).
    \item \textbf{/docs/img}: esta carpeta contiene las imágenes utilizadas en la documentación.
    \item \textbf{/docs/tex}: contiene los ficheros de \LaTeX de cada uno de los apartados de la documentación.
    \item \textbf{/src}: aquí encontramos los archivos \textit{.ipynb} con todo el código desarrollado. Dentro podemos encontrar el archivo \textit{Paquetes.ipynb} con los paquetes necesarios para ejecutar la aplicación, el archivo \textit{Pruebas.ipynb} con el entorno de pruebas y el archivo \textit{Turbina.ipynb}, desde el cual ejecutamos la aplicación.
    \item \textbf{/src/templates}: contiene el fichero \textit{.html} de la interfaz web.
\end{itemize}


\section{Manual del programador}

El objetivo del manual del programador es explicar el funcionamiento del proyecto a una persona que no haya participado en su realización y necesite conocer esta información. Este puede ser el caso de alguien que se haya contratado para el mantenimiento del sistema y necesite conocer de su funcionamiento.

\subsection{Entorno de desarrollo}

Para trabajar con el proyecto correctamente, es necesario instalar las siguiente serie de programas y dependencias:

\begin{itemize}
    \item \textbf{\href{https://www.python.org/downloads/}{Python}}: es el lenguaje en el que esta desarrollado casi todo el modelo. Es necesario para obtener todas las librerías para el correcto funcionamiento.
    \item \textbf{\href{https://jupyter.org/install}{Jupyter Notebook}}: es el IDE utilizado para el desarrollo del modelo. Es el recomendado para trabajar con el código, aunque se pueden usar otras aplicaciones.
     \item \textbf{\href{https://git-scm.com/book/en/v2/Getting-Started-Installing-Git}{Git}}: es necesario para hacer uso del repositorio de GitHub y poder obtener la última versión del proyecto.
\end{itemize}

\subsection{Descarga del código fuente}

Para obtener el código fuente y poder trabajar con él se deben seguir los siguientes pasos:

\begin{enumerate}
    \item Abrimos la terminal de Git Bash.
    \item Nos situamos, mediante el comando \textbf{cd}, en el directorio donde queramos guardar el proyecto.
    \item Utilizamos el comando \textbf{git clone} para clonar los documentos del repositorio en nuestro directorio. El comando completo para este caso es: \textbf{git clone https://github.com/Di3go99/Control-de-un-molino-eolico-usando-Aprendizaje-por-refuerzo-profundo.git}
    \item Una vez se complete la descarga, dispondremos de los archivos en nuestro directorio.
\end{enumerate}

Para trabajar con el código del proyecto, simplemente tendremos que abrir Jupyter Notebook y acceder al directorio donde hayamos guardado el proyecto. Tras esto, entramos en la carpeta \textbf{/src} y abrimos el archivo \textit{Turbina.ipynb}, donde encontraremos el código completo de la aplicación.


\section{Compilación, instalación y ejecución del proyecto}
\label{compilacion}

Lo primero que debemos hacer es instalar los paquetes necesarios para el código, los cuales podemos observar en la figura \ref{fig:paquetes}. Para ello, abrimos Jupyter Notebook, navegamos hasta la carpeta \textbf{/src} y abrimos el archivo \textit{Paquetes.ipynb}.

\imagen{Paquetes.png}{Paquetes necesarios para el modelo}{0.4}
\label{fig:paquetes}

Una vez cargado el archivo, pulsamos en el botón superior que dice \textbf{Run} y esperamos a que termine la ejecución de la celda.

Por último, volvemos a la carpeta \textbf{/src} y abrimos el archivo \textit{Turbina.ipynb}. Ahora hacemos clic en \textbf{Cell > Run All} (como puede verse en la figura \ref{fig:ejecucion}) para que se ejecuten todas las celdas de código.

\imagen{RunAll.png}{Botón para ejecutar todas las celdas}{0.8}
\label{fig:ejecucion}


\section{Pruebas del sistema}

Las pruebas realizadas sobre el proyecto, han sido pruebas del tipo ensayo-error realizadas durante todo el proceso de creación del mismo. A lo largo del transcurso del proyecto se han ido tomando los ajustes necesarios según el resultado de dichas pruebas.

En cuanto a las pruebas de calidad del código, se han realizado varios chequeos en Codacy\cite{Codacy} de cada una de las subidas al repositorio a lo largo del último mes. En la figura \ref{fig:codacy2} se puede ver un resumen con la calidad del código según Codacy

\imagen{Codacy.png}{Resumen de la calidad del código según Codacy}{1.1}
\label{fig:codacy2}
\apendice{Documentación de usuario}

\section{Introducción}

A continuación, se proceden a detallar, tanto los requisitos necesarios para que el usuario pueda usar el modelo, como los pasos para la instalación y el manual de uso del mismo.


\section{Requisitos de usuarios}

El usuario va a requerir del siguiente software.

\begin{itemize}
    \item \textbf{\href{https://www.python.org/downloads/}{Python}}: es el lenguaje en el que esta desarrollado casi todo el modelo. Es necesario para obtener todas las librerías para el correcto funcionamiento.
    \item \textbf{\href{https://jupyter.org/install}{Jupyter Notebook}}: es el IDE utilizado para el desarrollo del modelo. Es el recomendado para trabajar con el código, aunque se pueden usar otras aplicaciones.
    \item \textbf{\href{https://github.com/Di3go99/Control-de-un-molino-eolico-usando-Aprendizaje-por-refuerzo-profundo}{Última versión}}: es necesario que el usuario disponga de la última versión del proyecto descargada en su ordenador.
\end{itemize}


\section{Instalación}

Para la instalación se deben seguir los mismos pasos detallados en la sección \ref{compilacion} y, a continuación, abrir el enlace que podemos encontrar en el archivo \textit{Link.txt} del repositorio.


\section{Manual del usuario}

El modelo dispone de una interfaz muy sencilla, en la que el usuario simplemente debe rellenar un cuadro de texto y pulsar uno de los botones que aparecen en pantalla.

El funcionamiento es el siguiente:
\begin{itemize}
    \item \textbf{Cuadro de texto \textit{Potencia}}: el usuario debe rellenar, antes de pulsar ningún botón, este cuadro de texto con la cantidad de potencia que deseamos que el modelo alcance.
    \item \textbf{Botón \textit{Sin entrenar}}: mediante este botón, al usuario se le mostrará una gráfica del comportamiento del modelo previo al entrenamiento del agente.
    \item \textbf{Botón \textit{Entrenado}}: mediante este botón, al usuario se le mostrará una gráfica del comportamiento del modelo tras el entrenamiento del agente.
\end{itemize}


\bibliographystyle{plain}
\bibliography{bibliografiaAnexos}

\end{document}
